\documentclass[10pt,a4paper]{amsart}
\usepackage{amsmath,amssymb,amsthm,graphicx}
\usepackage{verbatim}
\newenvironment{reference}{\comment}{\endcomment}
%\newenvironment{reference}{}{}
\newenvironment{slogan}{\comment}{\endcomment}
\newenvironment{history}{\comment}{\endcomment}

% For commutative diagrams we use Xy-pic
\usepackage[all]{xy}

% We use 2cell for 2-commutative diagrams.
\xyoption{2cell}
% Theorem environments.
%
\theoremstyle{plain}
\newtheorem{theorem}[subsection]{Theorem}
\newtheorem{proposition}[subsection]{Proposition}
\newtheorem{lemma}[subsection]{Lemma}

\theoremstyle{definition}
\newtheorem{definition}[subsection]{Definition}
\newtheorem{example}[subsection]{Example}
\newtheorem{exercise}[subsection]{Exercise}
\newtheorem{situation}[subsection]{Situation}

\theoremstyle{remark}
\newtheorem{remark}[subsection]{Remark}
\newtheorem{remarks}[subsection]{Remarks}

\numberwithin{equation}{subsection}

% Macros
%
\def\lim{\mathop{\mathrm{lim}}\nolimits}
\def\colim{\mathop{\mathrm{colim}}\nolimits}
\def\Spec{\mathop{\mathrm{Spec}}}
\def\Hom{\mathop{\mathrm{Hom}}\nolimits}
\def\Ext{\mathop{\mathrm{Ext}}\nolimits}
\def\SheafHom{\mathop{\mathcal{H}\!\mathit{om}}\nolimits}
\def\SheafExt{\mathop{\mathcal{E}\!\mathit{xt}}\nolimits}
\def\Sch{\mathit{Sch}}
\def\Mor{\mathop{\mathrm{Mor}}\nolimits}
\def\Ob{\mathop{\mathrm{Ob}}\nolimits}
\def\Sh{\mathop{\mathit{Sh}}\nolimits}
\def\NL{\mathop{N\!L}\nolimits}
\def\CH{\mathop{\mathrm{CH}}\nolimits}
\def\proetale{{pro\text{-}\acute{e}tale}}
\def\etale{{\acute{e}tale}}
\def\QCoh{\mathit{QCoh}}
\def\Ker{\mathop{\mathrm{Ker}}}
\def\Im{\mathop{\mathrm{Im}}}
\def\Coker{\mathop{\mathrm{Coker}}}
\def\Coim{\mathop{\mathrm{Coim}}}

% Boxtimes
%
\DeclareMathSymbol{\boxtimes}{\mathbin}{AMSa}{"02}

%
% Macros for moduli stacks/spaces
%
\def\QCohstack{\mathcal{QC}\!\mathit{oh}}
\def\Cohstack{\mathcal{C}\!\mathit{oh}}
\def\Spacesstack{\mathcal{S}\!\mathit{paces}}
\def\Quotfunctor{\mathrm{Quot}}
\def\Hilbfunctor{\mathrm{Hilb}}
\def\Curvesstack{\mathcal{C}\!\mathit{urves}}
\def\Polarizedstack{\mathcal{P}\!\mathit{olarized}}
\def\Complexesstack{\mathcal{C}\!\mathit{omplexes}}
% \Pic is the operator that assigns to X its picard group, usage \Pic(X)
% \Picardstack_{X/B} denotes the Picard stack of X over B
% \Picardfunctor_{X/B} denotes the Picard functor of X over B
\def\Pic{\mathop{\mathrm{Pic}}\nolimits}
\def\Picardstack{\mathcal{P}\!\mathit{ic}}
\def\Picardfunctor{\mathrm{Pic}}
\def\Deformationcategory{\mathcal{D}\!\mathit{ef}}
\usepackage{cleveref}
\crefname{originalequation}{original equation}{original equations}
\Crefname{originalequation}{Original equation}{Original equations}
\makeatletter
\newlabel{equation-plain-versus-derived-algebras}{{0BZ4}{source}{}{sourcepointer-1}{}}
\newlabel{equation-plain-versus-derived-algebras@cref}{{[originalequation][0][]0BZ4}{[1][0][]source}{}{}{}}
\makeatother
\crefname{originalreference}{original reference}{original references}
\Crefname{originalreference}{Original reference}{Original references}
\makeatletter
\newlabel{lemma-compose-tensor-functors-general-algebra}{{0BZ5}{source}{}{sourcepointer-2}{}}
\newlabel{lemma-compose-tensor-functors-general-algebra@cref}{{[originalreference][0][]0BZ5}{[1][0][]source}{}{}{}}
\makeatother
\crefname{originallemma}{original lemma}{original lemmas}
\Crefname{originallemma}{Original lemma}{Original lemmas}
\makeatletter
\newlabel{more-algebra-lemma-K-flat-quasi-isomorphism}{{06Y0}{source}{}{sourcepointer-3}{}}
\newlabel{more-algebra-lemma-K-flat-quasi-isomorphism@cref}{{[originallemma][0][]06Y0}{[1][0][]source}{}{}{}}
\makeatother
\title{K-flatness makes the composition of derived tensor functors computed by plain tensors}
\author{The Stacks Project Authors}
\date{stacks-project:dga.tex:local-snapshot-20261012; GNU FDL 1.2 or later}
\begin{document}
\maketitle
\noindent\textit{Editorial note.} All delivered mathematical text is the author's own, copied from the pinned Stacks Project chapter file sources/stacks-dga/pinned-original-source.tex; no statement or proof step was rewritten, reordered or removed.\par\medskip
\noindent\textit{Editorial note.} The counted proof is the author's complete argument: dependencies carried as navigation stubs: the property (P) resolution construction behind equation equation-plain-versus-derived-algebras and Lemma lemma-compose-tensor-functors-general-algebra (both in-source), plus the standard More on Algebra K-flat quasi-isomorphism lemma (external tooling).\par\medskip
\noindent\textit{Editorial note.} The label lemma-compose-tensor-functors names the counted statement, and the proof environment immediately following it in the source is delivered verbatim as its proof.\par\medskip
\noindent\textit{Editorial note.} The counted statement is carried by the author's own lemma environment in the chapter source.\par\medskip
\noindent\textit{Allowed standard prerequisites (editorial).}\begin{itemize}
\item Property (P) resolutions compute derived tensor. Same unit, the construction behind equation equation-plain-versus-derived-algebras and Lemma lemma-compose-tensor-functors-general-algebra. The proof chooses a property (P) resolution P -> (A tensor\_R B)\_B and works with the plain tensor representatives of the derived functors as in that construction.
\item K-flat base change. More on Algebra, Lemma more-algebra-lemma-K-flat-quasi-isomorphism (standard tooling). Final step (source line 6090): tensoring a quasi-isomorphism with a K-flat complex of R-modules preserves quasi-isomorphism, so P tensor\_R C -> A tensor\_R B tensor\_R C is a quasi-isomorphism when C is K-flat.
\end{itemize}\medskip
\ifdefined\hypertarget\hypertarget{sourcepointer-1}{}\fi
\noindent\textit{Source reference.} Original equation 0BZ4: sources/stacks-dga/pinned-original-source.tex, source lines 5941--5941. Prerequisite pointer invoked at proof line(s) 6071; it names a standard fact of the same chapter that the counted argument uses, and every step of the counted proof built on it is retained verbatim in the delivered body.\par\medskip
\ifdefined\hypertarget\hypertarget{sourcepointer-2}{}\fi
\noindent\textit{Source reference.} Original reference 0BZ5: sources/stacks-dga/pinned-original-source.tex, source lines 5959--5959. Prerequisite pointer invoked at proof line(s) 6073; it names a standard fact of the same chapter that the counted argument uses, and every step of the counted proof built on it is retained verbatim in the delivered body.\par\medskip
\ifdefined\hypertarget\hypertarget{sourcepointer-3}{}\fi
\noindent\textit{Source reference.} Original lemma 06Y0: sources/stacks-chow/tags.txt, source lines 8869--8869. Prerequisite pointer invoked at proof line(s) 6090; it names a standard fact of a sibling chapter that the counted argument uses, and every step of the counted proof built on it is retained verbatim in the delivered body.\par\medskip
\begin{lemma}
\label{lemma-compose-tensor-functors}
Let $R$ be a ring. Let $(A, \text{d})$, $(B, \text{d})$, and
$(C, \text{d})$ be differential graded $R$-algebras.
If $C$ is K-flat as a complex of $R$-modules, then
(\ref{equation-plain-versus-derived-algebras})
is an isomorphism and the conclusion of
Lemma \ref{lemma-compose-tensor-functors-general-algebra} is valid.
\end{lemma}

\begin{proof}
Choose a quasi-isomorphism $P \to (A \otimes_R B)_B$ of differential
graded $B$-modules, where $P$ has property (P). Then we have to show
that
$$
P \otimes_B (B \otimes_R C) \longrightarrow
(A \otimes_R B) \otimes_B (B \otimes_R C)
$$
is a quasi-isomorphism. Equivalently we are looking at
$$
P \otimes_R C \longrightarrow
A \otimes_R B \otimes_R C
$$
This is a quasi-isomorphism if $C$ is K-flat as a complex of $R$-modules by
More on Algebra, Lemma \ref{more-algebra-lemma-K-flat-quasi-isomorphism}.
\end{proof}

\end{document}
